\section{Case 4: Sleep Schedule Breakdown}

\subsection*{Situation}
A person wants to sleep earlier but continues consuming stimulating content at
night until bedtime shifts later and later.

\subsection*{Framework map}
\begin{itemize}[nosep]
  \item \textbf{Current state:} high-stimulation late-night scrolling or media
        consumption.
  \item \textbf{Target state:} a designed pre-sleep regime that reliably
        transitions into sleep.
  \item \textbf{Dominant variables:} high $\Reward$, low resistance to
        continuation, weakly specified target state, and an environment that
        keeps brightness and stimulation active.
  \item \textbf{Missed decision window:} the first shutdown cue before the
        reward stream has fully captured the night.
\end{itemize}

\subsection*{Mechanism interpretation}
This is the case with the strongest direct connection to the research spine.
The sleep-transition literature gives a real example of patterned state change
that can be described in dynamical language. The manuscript should still stay
disciplined: the cited sleep work supports the importance of transition
structure and control parameters, not a full behavioural theory of personal
bedtime discipline. Even so, the case strongly benefits from thinking in terms
of regime change: if bright screens, novelty, and open-ended content keep the
system in the wrong basin, then ``just stop'' is not a real control strategy.

\subsection*{Evidence lane discipline}
This case has the strongest \textbf{strong} anchor available in the current
source set, because the indexed Nature Neuroscience sleep paper directly treats
falling asleep as a patterned transition dynamic. The translation from that
result to practical bedtime routine design remains partly \textbf{medium},
because the manuscript is still moving from biological regime transition to
everyday behavioural control design. That move is defensible, but it should be
stated as interpretation rather than finished proof.

\subsection*{Operational repair}
\begin{enumerate}[nosep]
  \item Replace the negative rule ``stop consuming'' with a positive pre-sleep
        state.
  \item Specify the object set: dim light, paper book, charger away from bed,
        and one fixed shutdown sequence.
  \item Treat the first shutdown cue as the key decision window rather than the
        final exhausted minute.
  \item Reduce parameter drift by keeping time, light, and device placement
        stable across evenings.
\end{enumerate}

\subsection*{Transfer lesson}
The more a target state is defined as the absence of a bad state, the weaker
its entry dynamics become. Productive control usually requires a replacement
regime, not only a prohibition.
